\emph{(i)} By Theorem \claimref{synthesis.positive-structure-is-twistor-form-at-imaginary-time}(i) the twistor form $\Sigma$ is $\mathrm{Spin}(2,4)$-invariant and the positive structure is $h(\psi, \varphi) = \Sigma(\psi, J\varphi)$ with $J = ie_{1}e_{+}$, $J^{2} = 1$ (in that theorem's labels, where $e_{1}$ and $e_{+}$ square to $+1$; in the labels of Proposition \claimref{synthesis.positive-conformal-energy-gives-positive-null-translations} they are $e_{0}$ and $\nu$, and $e_{1}e_{+}$ is $Z$). An element $g$ of $\mathrm{Spin}(2,4)$ preserves $\Sigma$, so it preserves $h$ if and only if $\Sigma(g\psi, Jg\varphi) = \Sigma(g\psi, gJ\varphi)$ for all $\psi, \varphi$, that is, by nondegeneracy, $gJ = Jg$: the symmetry group of $h$ is the centraliser of $Z$. By (iii) of that theorem every positive structure is $h_{m}$ for a forward-tube point $m$, and the group acts transitively on the tube (the bounded symmetric domain of $SU(2,2)$, cited from memory), so for every positive structure the symmetry is a conjugate of this centraliser, and the statements below hold for it after conjugation. In $\Cl(2,4)$ a bivector $B$ fixes a vector $v$ under the adjoint action of $\exp(tB)$ for all $t$ if and only if $Bv - vB = 0$. The check computes, exactly, the rank of $v \mapsto [Z, v]$ on the six-dimensional vector space (rank $2$, kernel $\mathrm{span}(e_{1}, e_{2}, e_{3}, e_{-})$), of the joint map for the six bivectors $e_{i}e_{j}$ of that span (kernel $\mathrm{span}(e_{0}, \nu)$), and of the joint map for all seven (kernel $0$), and that the six commute with $Z$, so the seven span the centraliser's Lie algebra $\mathfrak{u}(1)\oplus\so(4)$ (the centraliser of $Z$ in $\so(2,4)$ has dimension $1 + 6$: $Z$ itself and $\Lambda^{2}$ of its orthogonal complement; the mixed bivectors $e_{0}e_{i}$, $\nu e_{i}$ do not commute with $Z$). Controls: $e_{1}e_{2}$ alone fixes a four-dimensional space, and $Z + e_{0}e_{1}$ does not commute with the $\so(4)$. Signatures: $e_{0}^{2} = \nu^{2} = 1$ and the other four square to $-1$, so $Q > 0$ on the plane and $Q < 0$ on its complement; with $Q(I) = \Lambda/3$ the fixed space of $\mathfrak{u}(1)$ consists of $\Lambda < 0$ tractors and that of $\so(4)$ of $\Lambda > 0$ ones. $\nInf = \nu + e_{-}$ is null and has components in both, so it lies in neither. The rotor $R = 1 + tZ$ has $R\tilde R = 1 + t^{2}$ and maps $\nu$ to $(1 + t^{2})$ times a unit vector of the plane, with nonzero $e_{0}$ component (the Cayley parametrisation reaches every rational point of the circle), checked at $t = \tfrac12, 3, -\tfrac27$. Orbits: the centraliser acts on $\mathrm{span}(e_{0}, \nu)$ through $SO(2)$ and on the complement through $SO(4)$, independently, so the orbit of $I = I_{+} + I_{-}$ is the product of the circle of radius $Q(I_{+})^{1/2}$ and the $3$-sphere of $\{Q = Q(I_{-})\}$ in the complement; it is a point only for $I = 0$. A principle invariant under the positive structure therefore selects at most a union of orbits, through the two invariants. With the boost $K = e_{0}e_{1}$ added, the joint kernel of $[K, \cdot]$ and the $\so(4)$ commutators is one-dimensional, spanned by $\nu$ (checked). For the tube point $m = i\rho e_{0}$, $X(m) = \nO + \tfrac{\rho^{2}}{2}\nInf + i\rho e_{0}$; its real and imaginary parts span the positive plane of $h_{m}$ (for $\rho = 1$ this is $\mathrm{span}(\nu, e_{0})$, the plane of $Z$; in general the plane of the conjugate of $Z$, argued from the transitivity above), $Q(\mathrm{Re}\,X) = \rho^{2}$, and $\tfrac{2}{\rho^{2}}\mathrm{Re}\,X = \nInf + \tfrac{2}{\rho^{2}}\nO$, the de Sitter tractor at $\Lambda = 12/\rho^{2}$, which commutes with $K$; the check verifies this at four values of $\rho$. This is the exact form of the matching $\rho^{2} = 12/\Lambda$ between the frozen point and the boost-compatible tube points of synthesis.wick-rotated-boost-selects-observer-tube-point (filed alongside): that result selects the tube point from a boost and an observer, this one shows the tube point alone does not select the scale. In the de Sitter normalisation $\sigma = 1 - \Lambda Q/12$, $\ell = (3/\Lambda)^{1/2}$, this is $\rho = 2\ell$, a factor $2$ from the argued normalisation $\rho = 1/\kappa = \ell$ there; both are stated only for the foot $x_{\perp} = 0$. The tabulated scales are $\ip{X(x)}{\nu} = \tfrac12(1 - Q(x))$, $\ip{X(x)}{e_{0}} = x^{0}$, $\ip{X(x)}{e_{-}} = \tfrac12(1 + Q(x))$, checked at two points; $\tfrac12(\nInf + 2\nO) = \nu$ identifies $\nu$ with the de Sitter tractor at $\Lambda = 12$ up to the factor.

\emph{(ii)} For $g = \varphi^{2}\eta = e^{2w}\eta$, $w = \ln\varphi$, the Christoffel symbols are $\Gamma^{c}{}_{ab} = \delta^{c}_{a}\partial_{b}w + \delta^{c}_{b}\partial_{a}w - \eta_{ab}\eta^{cd}\partial_{d}w$, and $\partial w = \partial\varphi/\varphi$, $\partial\partial w = \partial\partial\varphi/\varphi - \partial\varphi\,\partial\varphi/\varphi^{2}$ are rational in the jets of a polynomial $\varphi$. The check computes $\Ric_{ab} = \partial_{c}\Gamma^{c}{}_{ab} - \partial_{b}\Gamma^{c}{}_{ac} + \Gamma^{c}{}_{cd}\Gamma^{d}{}_{ab} - \Gamma^{c}{}_{bd}\Gamma^{d}{}_{ac}$ exactly at three rational points (the origin and two generic points) for the two counter-models and a control, and verifies $\varphi^{4}\mathcal{R} = -6\varphi\Box\varphi$ in each case ($\mathcal{R} = g^{ab}\Ric_{ab}$; the overall sign is the one of this index convention, the store's convention flips it, which changes nothing below). In the action, $\sqrt{|g|} = \varphi^{4}$ and $-6\varphi\Box\varphi = 6\,\eta(\partial\varphi, \partial\varphi)$ up to a divergence, so the restricted action is $\int(6\,\eta(\partial\varphi, \partial\varphi) - 2\Lambda\varphi^{4})$ and its variation in $\varphi$ gives $-12\Box\varphi - 8\Lambda\varphi^{3} = 0$, equivalently $\mathcal{R} = 4\Lambda$ by the identity. For $\varphi = 1 + x_{1}x_{2}$ and $\varphi = 1 + \tfrac13 t + x_{1}^{2}x_{3} - x_{2}^{2}x_{3}$ the check finds $\Box\varphi = 0$ as a polynomial identity, hence $\mathcal{R} = 0$ everywhere and both are stationary at $\Lambda = 0$, while $\Ric_{ab} - \tfrac14\mathcal{R}g_{ab} \ne 0$ at every sample point; at the origin for the first, $\partial w = 0$ and the trace-free Ricci tensor is $-2\,\partial_{1}\partial_{2}w = -2$ in the $(1,2)$ slot. A scale is almost-Einstein exactly when the trace-free Ricci tensor of its metric vanishes (the derived-metric reading of the parallel scale tractor, Fact \claimref{correspondence.imported-vacuum-equations-scale-tractor}, Fact \claimref{synthesis.imported-gtg-dictionary}), so these stationary scales are not almost-Einstein. Control: $\varphi = 1 + x_{1}^{2}$ has $\Box\varphi = -2 \ne 0$ and $\mathcal{R} \ne 0$, with the identity still holding; $\varphi = 1$ is flat. The cross-subject reading: this is the Yamabe problem, in which the conformal factor of a metric carries only the scalar curvature (cited from memory).

\emph{(iii)} In the flat splitting, with $\Sch = 0$ and $J = 0$ for $\eta$, the tractor of the scale is $D\sigma = (\sigma, \nabla\sigma, -\tfrac14\Box\sigma)$ and the flat tractor connection $\nabla_{a}(\sigma, \mu, \rho) = (\partial_{a}\sigma - \mu_{a}, \partial_{a}\mu_{b} + \eta_{ab}\rho, \partial_{a}\rho)$ (Fact \claimref{correspondence.imported-vacuum-equations-scale-tractor}; Theorem \claimref{synthesis.infinity-vacuum-incidence-is-a-property-of-the-scale}(4), whose middle slot is $-\Sch^{\circ}$, agrees at points where $\sigma = 1$ and $d\sigma = 0$) gives $\nabla_{a}D\sigma = (0, H^{\circ}_{ab}, -\tfrac14\partial_{a}\Box\sigma)$. With the tractor metric $\ip{(\sigma, \mu, \rho)}{(\sigma', \mu', \rho')} = \sigma\rho' + \rho\sigma' + \eta(\mu, \mu')$ the zero top slot kills the cross terms, so $\eta^{ab}\ip{\nabla_{a}D\sigma}{\nabla_{b}D\sigma} = H^{\circ}\!:\!H^{\circ}$. Varying, $\delta\int H^{\circ}\!:\!H^{\circ} = 2\int H^{\circ ab}\partial_{a}\partial_{b}\delta\sigma$ (the trace part drops against the trace-free $H^{\circ}$), so the Euler-Lagrange equation is $\partial_{a}\partial_{b}H^{\circ ab} = \Box^{2}\sigma - \tfrac14\Box^{2}\sigma = \tfrac34\Box^{2}\sigma = 0$ up to the factor $2$; the check verifies the operator identity $\partial_{a}\partial_{b}H^{\circ ab} = \tfrac34\Box^{2}\sigma$ on a polynomial of degree four with nonzero $\Box^{2}$. For $\sigma = 1 + x_{1}x_{2}$ the operator vanishes and $H^{\circ}$, whose only entry is the $(1,2)$ one, is nonzero at every sample point. For $\sigma = (t + x_{1})^{2}$, $\Box^{2}\sigma = 0$, and $H^{\circ} = H = 2\,l^{\flat}\otimes l^{\flat}$ with $l = e_{0} + e_{1}$, so $H^{\circ}\!:\!H^{\circ} = 4\,\eta(l, l)^{2} = 0$ with $H^{\circ} \ne 0$ (checked at three points). Control: the almost-Einstein scale $1 - \tfrac12 Q(x)$ has $H^{\circ} = 0$, and changing one coefficient gives $H^{\circ} \ne 0$.

\emph{(iv)} (Argued, not proved.) A local functional $S[\sigma]$ has the single Euler-Lagrange expression $\delta S/\delta\sigma$. The almost-Einstein condition is $H^{\circ} = 0$ in the flat class, nine independent second-order conditions, whose solutions are the six-parameter family $c + b\cdot x - \tfrac{k}{2}Q(x)$ (stored). If $\delta S/\delta\sigma = 0$ is a scalar equation of order $k$ that is noncharacteristic on some hypersurface, the Cauchy-Kovalevskaya theorem gives a solution for every analytic set of $k$ Cauchy data on it, an infinite-dimensional family, so its solution set is not the six-dimensional family; the degenerate (everywhere characteristic) case is not treated. A principle with exactly the almost-Einstein stationary points therefore needs more than one equation per scale: Lagrange multipliers $B^{a}$ with a term $\int\ip{B^{a}}{\nabla_{a}I'}$, whose equations restate $\nabla I' = 0$, or a variation of the conformal class as well, which is a dynamical principle for the gauge field $\hb$ (whose Euler-Lagrange equations for the Einstein-Hilbert action are the full field equations, the trace-free part included). The open problem excludes the first. With matter, the second is the field equations with a stress tensor, which Theorem \claimref{synthesis.infinity-vacuum-incidence-is-a-property-of-the-scale}(4) and Theorem \claimref{synthesis.vacuum-field-stress-tensor-is-einstein-tensor-or-trace} already place outside the vacuum field's data.
